\section{优化技巧：按作用对象分四类}
\label{sec:technique}

前一节讲的是「分不分、怎么分」。但一个已经上了分离的部署，
真正决定它跑得好不好的往往不是架构，而是\textbf{一串具体技巧}。
本节把这些技巧组织起来。

组织的维度不是「属于哪个框架」，也不是「新不新」，而是
\textbf{作用于哪个环节}。按此分四类\texttt{[c-technique]}：

\begin{table}[t]
\centering
\caption{优化技巧按作用对象分四类。同类的技巧互相替代，
异类可以叠加——故取舍问题不是「选哪个」，而是「每类选一个」。}
\label{tab:technique}
\small
\begin{tabular}{@{}p{1.15cm}p{2.0cm}p{2.5cm}p{1.3cm}@{}}
\toprule
类别 & 作用于 & 代表技巧 & 实测幅度 \\
\midrule
A 省传输量 & KV 字节数本身 & 逐层传输、chunk 级传输、
  KV 压缩/量化、前缀缓存 & 成比例 \\
B 让传输不可见 & 传输与计算的时序 & TBO、SBO、DBO、piggyback
  & 隐藏传输 \\
C 均衡负载 & 每个 rank 分到多少 & Waterfill、LPLB、EPLB、
  请求迁移、弹性配比 & 0.84--7.34\% \\
D 省 decode 步 & 要跑多少步 & DSpark 自适应验证、投机解码
  & 前沿 \\
\bottomrule
\end{tabular}
\end{table}

四类的性质完全不同，混着谈会得出错误结论。以下逐类说。

\subsection{A 类：省传输量}

最直接的一类——传输成本与字节数成正比，那就少传。

\noindent\textbf{逐层传输。} 不等整个 prefill 算完再传，而是每算完
一层就异步传该层的 KV，prompt 计算同时进入下一层。Splitwise 的做法是：
在每层末尾触发该层 KV 的异步传输，让 prompt 计算继续
\texttt{[c-technique]}。好处有三——传输被计算掩盖、token 阶段可以
更早开始、prompt 机器可以更早释放显存。

\noindent 但这条技巧\textbf{有适用下限}，这点常被忽略：原文明确指出，
小 prompt 时 KV 很小，「没有必要为逐层传输所需的细粒度同步付出开销」
\texttt{[c-technique]}。也就是说，逐层传输把一次大传输切成 $L$ 次
小传输，每次都付固定开销；当 KV 总量小到固定开销占比显著时，
切得越细越亏。

\noindent\textbf{chunk 级传输。} 与逐层同类但切分维度不同：沿
chunked prefill 的 chunk 切，使 prefill 与 KV 传输并行。
另有 request 级传输，用来减少传输次数（一次传整个请求的 KV）。
原文把三者的关系说清了：chunk 级让计算与网络并行，request 级
最小化传输次数，逐层则可与 chunk 级结合以进一步并行
\texttt{[c-technique]}。

\noindent 这三者是同一维度的不同粒度。选哪个取决于两件事：
固定开销占多大（粒度太细不划算）、以及 prefill chunk 有多长
（chunk 级需要 chunked prefill 才有意义）。

\noindent\textbf{KV 压缩与量化。} 直接从字节数下手。第
\ref{sec:model} 节已提到逐阶段量化的代价不对称：只量化 prefill 的
精度损失是只量化 decode 的 $1.1$ 至 $4.1$ 倍\texttt{[c-model]}。
所以压缩要用在 decode 侧更划算——那里 KV 读得多、且对精度更宽容。

\noindent\textbf{前缀缓存。} 相同前缀不重算、也不重传。
严格说它省的是 prefill 计算而非传输，但它会\textbf{改变传输量的分布}：
命中前缀后需要搬的 KV 变小，于是前面所有按「全量 KV」算的成本估计
都要下调。这是部署时最容易低估的一条。

\noindent\textbf{重算 vs 传输。} 还有一个非技巧的技巧：KV 太小时，
\textbf{重算比传输便宜}。这条与第 \ref{sec:boundary} 节的
「上下文极短导致固定开销主导」是同一件事的两面。

\subsection{B 类：让传输不可见}

不减少字节数，而是让传输发生在计算的同时，从而在端到端时间里
看不见它。这一类全部依赖\textbf{MoE 的 all-to-all 通信}与
\textbf{attention 计算}可以交错这一事实。

\noindent\textbf{TBO（Two-Batch Overlap）。} 把请求拆成两个
micro-batch，让一个 micro-batch 的 attention 计算与另一个的
dispatch/combine 通信交错。SGLang 的官方描述是：TBO 把请求拆成
micro-batch，把 attention 计算与 dispatch/combine 操作交错，
执行图中的 yield 点允许暂停以形成重叠，从而在\textbf{不增加峰值显存}
的前提下提高吞吐\texttt{[c-technique]}。

\noindent 最后半句是关键：TBO 不靠加显存换吞吐，靠的是给通信
找一个可以并行的伙伴。

\noindent\textbf{SBO（Single-Batch Overlap）。} 不拆 batch，
而是在同一个 micro-batch 内部把通信与 shared expert 的计算重叠
\texttt{[c-technique]}。相比 TBO，SBO 不需要拆 batch，
故没有「拆完每半太小、GEMM 效率下降」的风险，但可重叠的窗口也更窄
（只有 shared expert 那一段）。

\noindent\textbf{DBO（Dual Batch Overlap）。} vLLM 的实现思路相同但
手段不同：在 model runner 里把 batch 拆成两个 micro-batch，
\textbf{起两个 worker 线程}各跑一半，在两个线程之间设置 yield 点，
使一个在算时另一个在等通信，形成乒乓\texttt{[c-technique]}。

\noindent DBO 有一段值得记的工程细节：拆不拆要\textbf{跨所有 DP rank
统一决定}——只要有一个 rank 拆不动，全体都不拆；且若 padding 后
某个 rank 的第二个 micro-batch 为空，也全体放弃\texttt{[c-technique]}。
这是「并行组内必须同步决策」的典型代价。

\noindent 更实用的是它的\textbf{阈值}：DBO 只在 batch 的 token 数超过
阈值时才启用，低于阈值不拆\texttt{[c-technique]}。原因不难推：
拆成两半后每半的 GEMM 变窄、线程切换有成本，小 batch 上得不偿失。

\noindent\textbf{piggyback。} 让 decode 搭 prefill 的便车一起批处理
（Sarathi-Serve 的做法：把 decode 挂到 chunked prefill 上）。
它在不分离的架构里替代分离——这一点很重要：
\textbf{B 类技巧里有几个是分离的替代品，不是补充}。

\noindent 这一类的共同前提是 \textbf{MoE + EP}。若模型是 dense，
没有 all-to-all 可重叠，B 类基本失效。这是选型时第一个要看的。

\subsection{C 类：均衡负载}

MoE 的 router 不产生均匀的专家流量，某些专家收到的 token 远多于
其它，EP 组于是等最忙的那个 rank。C 类技巧处理这件事。

\noindent\textbf{Waterfill。} 一个轻量的 shared-expert 负载均衡：
把 shared expert 通过 DeepEP 派发，并\textbf{指派给负载较轻的 rank}
\texttt{[c-technique]}。实测：在两台 Hopper 节点上跑 DeepSeek-V3/R1
风格负载，Waterfill 在 MMLU、GPQA、GSM8K 上把总吞吐提升
\textbf{+1.48\% 至 +4.66\%}；在 DeepSeek V4 上最好测点从
49{,}253\,tok/s 提到 51{,}677\,tok/s（+4.92\%）\texttt{[c-technique]}。

\noindent 注意幅度：个位数百分比。这类技巧\textbf{不是数量级的改进}，
但它的成本也很低（只改派发目标，不动权重、不改路由）。

\noindent\textbf{LPLB。} 针对\textbf{冗余专家副本}的线性规划均衡器：
对每层解一个 dispatch 优化问题，在冗余副本之间重新分配 token
\texttt{[c-technique]}。实测：同样的两台 Hopper 节点、同样的
冗余 EPLB 放置，LPLB 把总吞吐提升 \textbf{+0.84\% 至 +7.34\%}
\texttt{[c-technique]}。

\noindent LPLB 有一个必须说清的性质：它\textbf{不移动权重，也不改变
router 的逻辑决策}，只改物理副本的分配\texttt{[c-technique]}。
故它不损精度——这一点与量化类技巧形成对比，也是它能「默认打开」
的原因。

\noindent\textbf{两个前置依赖，缺了就只有开销。}

第一，\textbf{必须有冗余副本}。原文明确：在 \texttt{red0}
（无冗余副本）配置下，没有副本可重新平衡，算法无法改善派发，
其 all-reduce 与求解路径\textbf{纯粹表现为开销}\texttt{[c-technique]}。

第二，\textbf{收益取决于在线流量的分布与标定分布的偏离程度}。
原文的说法是 LPLB 的收益「跟随 live batch 偏离标定分布的程度」，
并在中等规模处最强\texttt{[c-technique]}。

\noindent 这两条合起来说明一件事：C 类技巧的收益是
\textbf{有条件的、个位数百分比},不该被当成「开了就变快」的开关。

\noindent\textbf{PD 给 C 类带来的新约束。} 这是本节最需要单独讲的
一处交互。DP attention 下，\textbf{不同的 EP rank 在同一步跑不同的
forward 模式}——有的在跑 prefill、有的在跑 decode、有的 idle。
于是没有单个 rank 能看到全局的 token 分布\texttt{[c-technique]}。

\noindent 后果很具体：均衡算法的输入（全局各专家的 token 计数）
在 PD 下不是本地可得的，必须额外做一次 all-reduce 才能拼出来。
这既是开销，也是延迟——负载均衡决策要等全局信息到齐。

\noindent 这一条解释了为什么有些在不分离部署里跑得很好的均衡技巧，
搬到 PD 上收益变薄：\textbf{信息可见性被架构改变了}。

\noindent\textbf{其它 C 类手段。} EPLB（静态专家放置）、
请求迁移（Llumnix 在运行时重调度请求）、以及弹性配比
（FluidPD 在运行时改 P:D 比）。它们作用于不同的时间尺度：
EPLB 是天级、请求迁移是秒级、弹性配比是分钟级。
\textbf{不同时间尺度的手段可以叠加}，因为它们不抢同一个决策点。

\subsection{D 类：省 decode 步}

decode 的每一步都要读全部权重与 KV，故减少步数是最直接的加速。
投机解码就是这条路：用小模型起草 $k$ 个 token，大模型一次验证。

\noindent 但这条路有一个\textbf{随负载反转}的收益结构。原文说得很清楚：
batch size 为 1 时这是笔好交易——GPU 受显存带宽限制、算力有余，
起草的额外计算近乎免费；到 batch size 256 时这笔交易就微妙了，
draft token 开始与真 token 抢同一份算力，每个被拒的 token 都浪费
了有用的算力，被拒够多时吞吐显著下降\texttt{[c-technique]}。

\noindent 为什么会拒？因为\textbf{逐位置存活率衰减很快}：在
DeepSeek-V4-Pro-0813 上，一个 7 token 块的最后一个 token 存活率
不足 10\%，而第一个超过 70\%\texttt{[c-technique]}。那个低概率的
token 在每次验证批里都占一个槽位——显存受限时这个槽位是免费的，
算力饱和后它就有真实成本。

\noindent\textbf{DSpark 的做法}：给每个 draft token 打一个置信度，
\textbf{按步决定验证多少}，而不是按部署选一个固定的
\texttt{num\_speculative\_tokens}\texttt{[c-technique]}。

\noindent 效果：开启自适应验证后，投机解码在并发 1 到 256 全程都能
带来收益，且低并发时仍保有长 draft 的好处\texttt{[c-technique]}。
原文的说法是「停留在 Pareto 前沿的边缘」\texttt{[c-technique]}。

\noindent 这条技巧的意义超出它自身：它是\textbf{把静态超参换成按负载
决策}的一个范例。前面 A、B、C 三类里的阈值（DBO 的 token 阈值、
LPLB 的冗余度）都可以用同样的思路做成自适应。

\noindent\textbf{代价与限制。} DSpark 需要 ALWAYS 级的 CUDA graph
支持，故不支持 \texttt{--enforce-eager}、LoRA 与 pipeline
parallel\texttt{[c-technique]}。这是 B、D 两类共有的特征：
\textbf{越依赖 kernel 级重叠的技巧，对执行模式的要求越硬}。

\subsection{补充三类：数据通路、准入、并行度重切分}

上面四类是按「作用于哪个环节」分的。还有三类技巧不落在这四类里，
但都是 PD 部署上实际在用的，且各有独立的量级。

\noindent\textbf{\sym{①} 数据通路层：不是少传，而是让每次传输更便宜。}

最典型的是 \textbf{zero-copy}：让发送/接收 kernel 直接操作应用的数据区，
省掉中间那一次拷贝\texttt{[c-technique]}。在 CloudMatrix384 的
send/recv 设计里，除了基于内存语义（MTE2/MTE3）的传输路径，
还专门提供 zero-copy 版本\texttt{[c-technique]}。

\noindent 与之配套的是 \textbf{GPUDirect RDMA}：让网卡直接读写显存，
绕过主机内存。Mooncake 把 KV 块在 CPU 与 GPU 之间的搬移交给一个
独立的 GPUDirect RDMA 组件（Messenger）\texttt{[c-technique]}。

\noindent 这一类的收益不在「传得少」，而在\textbf{把传输从「要经过
CPU 与两次拷贝」变成「网卡直接进显存」}。它与 A 类正交，可以叠加。

\noindent\textbf{\sym{②} 准入与拒绝：容量不够时不排队，直接拒。}

这条在过载场景下最重要，也最常被忽略。传统研究假设所有请求都会被处理，
但真实的高负载场景下，\textbf{把注定超时的请求排进队列只会拖垮其它请求}。

\noindent Mooncake 的做法是\textbf{预测式 early rejection}：先预测这条
请求能否在 SLO 内完成，不能就早拒\texttt{[c-technique]}。效果：在
某些模拟场景下吞吐最高提升 \textbf{525\%} 且仍满足 SLO；在真实负载下
让 Kimi 多处理 \textbf{75\%} 的请求\texttt{[c-technique]}。

\noindent 这是本文所引技巧里\textbf{幅度最大的一条}，但它提升的不是
「单位时间算得更多」，而是「单位时间算的都是有用的」。二者在过载时
差别巨大。

\noindent 这也解释了为什么它属于 PD 的优化技巧：分离之后 P 与 D 的
队列是独立的，一条请求要过两个池子，\textbf{任一池子超时都白做}，
故准入判断必须放在分离之前，且要同时考虑两段的排队时间。

\noindent\textbf{\sym{③} 并行度重切分：运行中改 DP/TP/PP 布局。}

比「弹性改 P:D 比」更激进的一层——不只改两阶段的比例，而是改
\textbf{每个阶段内部的并行方式}，并为此重新切分权重与重算 KV。

\noindent 有一个完整的实测：某系统以 Layout A 起步，当活跃轨迹数降到
批量的 10\% 以下时切到 Layout B（DP=8、PP=16、TP=8），
切换要付一次\textbf{权重 reshard 与 KV 重物化}的代价
\texttt{[c-technique]}。收益：rollout makespan 从 528.8\,s 降到
259.1\,s，有效 decode 吞吐提升 2.04 倍\texttt{[c-technique]}。

\noindent 原文把这归因为「相位匹配的并行度」：Layout A 吸收密集突发，
而重新切分让稀疏尾部的 decode 效率回升\texttt{[c-technique]}。

\noindent 这一类的取舍很清晰：\textbf{重切分有一次性成本，
故只在相位持续足够久时才划算}。这与第 \ref{sec:boundary} 节的
「窗口」是同一个道理——收益要摊得平固定成本。

\subsection{怎么选：三条判据}

把四类放在一起，取舍可以按三条判据走。

\noindent\textbf{\sym{①} 先看模型是不是 MoE + EP。} 若不是，B 类（TBO/SBO/
DBO）基本不可用，因为没有被重叠的 all-to-all。这一步就能砍掉一半选项。

\noindent\textbf{\sym{②} 每类选一个，先做前置条件检查。} 同类别替代，
故不必都上。但每条都有前置条件：逐层传输要求 KV 够大（否则固定开销
主导）；LPLB 要求有冗余副本（\texttt{red0} 下只有开销）；
DBO 要求 token 数过阈值；DSpark 要求 CUDA graph 全量模式。
\textbf{前置条件不满足时，开了是负收益}——这是本文最想强调的一点。

\noindent\textbf{\sym{③} 按幅度排优先级。} A 类与 C 类不在一个量级：
A 类按传输比例省（可能省掉一半），C 类是个位数百分比
（$0.84\%$--$7.34\%$）。D 类的幅度取决于并发——低并发时大、
高并发时靠自适应才不亏。

\noindent 所以合理的顺序是：\textbf{先用 A 类把量降下来，
再用 B 类把剩下的藏起来，然后用 D 类减少步数，
最后用 C 类抹平尾部不均}。反过来做（先上 C 类）会得到
「调了很久只有 3\%」的体验。

\noindent 最后一条经验值得单独记：\textbf{这四类里有几条是分离的替代品
而非补充}。piggyback（decode 搭 prefill 的便车）与 chunked prefill
在不分离架构里就能拿到相当一部分收益。若一个部署还没上分离，
先试 A 类与 piggyback 的性价比通常高于直接上分离——这与第
\ref{sec:boundary} 节的判据一致：分离是这些手段里最贵的一个。
